\scanpage{11}
tout $s\in S$, $P_s$ le \unsure{torseur|torseurs} sur $G(\unsure{\mathbb{Q}|Q})$
\[
 P_s=\operatorname{Isom}_{\otimes}(\unsure{\omega|w}\inserted{\text{foncteur oubli}},\Phi_sF)
\]
Pour $s$ variable, les $P_s$ forment un revêtement principal \inserted{(à droite)} $P$ de $S$, de groupe $\Gamma=G(\mathbb{Q})$. On a tout ce qu'il fallait pour que le \unsure{$\otimes$} composé
\diagram{\begin{tikzcd}[ampersand replacement=\&,column sep=large]
 C=\operatorname{Rep}(G) \arrow[r,harpoon,"F"] \& \operatorname{Hodge}(S) \arrow[r,harpoon,"{V\mathrel{\mapstochar\rightharpoonup} V\times_S P}"] \& \operatorname{Hodge}(P)
\end{tikzcd}}
\deleted{\unsure{ait une struct|ait un struct}} definisse sur $P$ une $G$-subordination, d'où un morphisme
\[
 g:P\rightharpoonup\unsure{\mathscr{J}|\mathcal{J}|\gamma}.
\]
\deleted{Ceci posé,} Ce dernier est compatible avec l'action de $\Gamma=G(\mathbb{Q})$, \inserted{(opérant sur $\mathscr{J}$ à \uline{droite}).} Et la catégorie des \deleted{\unsure{$\otimes$ foncteurs} $F$} $G$-subordinations $F$ sur $S$ compatibles avec le foncteur oubli est équivalente à celle des couples $(P,g)$, $P$ un rév. principal de $S$ de gpe $\Gamma$, $g:P\relbar\mathscr{J}$ un $\Gamma$ morphisme.

Soit $\Gamma_0\subset\Gamma=G(\mathbb{Q})$ un ss-gpe [de \unsure{préférence|préference|preférence}, commensurable à un $G(\mathbb{Z})$]. On appelle $G$-subordination $\Gamma_0$-\unsure{rigidifiée|rigidifiee} sur $S$ la donnée d'une rev. principal $P_0$ de $S$ de
% archive sheet number: 10
